\chapter*{卷四\quad 后题}
\phantomsection
\addcontentsline{toc}{chapter}{卷四\quad 后题}
\markboth{卷四\quad 后题}{卷四\quad 后题}

\noindent
卷四六片，都不是当代之物。它们或为后世题记，或为伪托，或为抄手之讹。

此卷入书，是因为它们都在原件上。一片上的字若非当时所刻，其史料价值反而更大：它说明此后仍有别人来过这里。

---

\begin{jian}{简三十二}
后三百七十二年，有客至。门不可启，壁上题名而去。
\end{jian}
\chuchu{原件残片 81 号，长 13.4 cm，金属，字口浅，与他片不类。}
\xiaojì{「后三百七十二年」若自前 221 年计，则为公元 151 年。此时中原为东汉。此题记之人自称为「客」，不称职衔，亦不留名。}

\begin{jian}{简三十三}
此非月也，此舟也。舟已不行。吾辈守之三世矣。
\end{jian}
\chuchu{原件残片 83 号，长 18.1 cm，金属。}
\xiaojì{「此舟也」一语，与《战争的雷霆》全书之核心判断相同。然此片字迹拙涩，刻痕不均，疑为后世摹刻，非原件。详下编第七章。}

\begin{jian}{简三十四}
大明万历十七年，内官监奉旨采办，至此而返。
\end{jian}
\chuchu{原件残片 85 号，长 11.2 cm，陶瓷，背面有朱印一方，不可辨。}
\xiaojì{「大明万历十七年」，即一五八九年。此为全书最晚的一条自记年代。所谓「至此而返」，此字何指，无从判定；若指广寒宫，则明代有人至月之说，始见于此。}

\begin{jian}{简三十五}
□有十二，不可尽。尽者三，余九亡。
\end{jian}
\chuchu{原件残片 87 号，长 9.8 cm，金属，上下俱蚀。}
\xiaojì{「十有二」与玉山十二个乘员舱之数合，然文中无一处可与舱对上。「尽者三，余九亡」，三与九之和为十二，此数可注意。}

\begin{jian}{简三十六}
一物，圆，径不及二寸，重若干。非金非石。击之不鸣。来源不明，暂不归类。
\end{jian}
\chuchu{原件残片 89 号，长 12.6 cm，金属。此片所述之圆片，现存基金会，编号 ФВ-1963-46。}
\xiaojì{「非金非石」「击之不鸣」，形容与《战争的雷霆》译者序所称之圆片合。该器物经俄方分析，铼锇同位素异常，地球无对照矿源。「暂不归类」四字为一九六三年编目时之批注，非原件文字。}

\begin{jian}{简三十七}
字尽于此。
\end{jian}
\chuchu{原件残片 91 号，长 6.2 cm，金属，为一残片之末段，前文俱失。}
\xiaojì{「字尽于此」可为抄手自记，亦可为全书之末。以本书编次论，取后义。}
